\uuid{hsWK}
\titre{Primitives usuelles}
\auteur{hanani ; mzari}
\organisation{exo7}
\licence{CC BY-NC-SA 4.0}
\datecreate{2020-01-20}
\niveau{L1}
\module{Analyse}
\chapitre{Calcul d'intégrales}
\sousChapitre{Primitives diverses}
\difficulte{1}
\theme{Primitives, primitives usuelles}
\exo7qcmid{lille-2.primitives.007}

\contenu{
\begin{question}
\qclassification{127.11}
\qtags{motcle=primitives usuelles}
Parmi les affirmations suivantes, cocher celles qui sont vraies. Sur $]-1,+\infty[$, on a :
\begin{answers}
    \bad{$\displaystyle \int \frac{\mathrm{d}x}{(x+1)^3}=\frac{-3}{(x+1)^4}+k$, où $k\in \Rr$.}
    \good{$\displaystyle \int \frac{\mathrm{d}x}{(x+1)^3}=\frac{-1}{2(x+1)^2}+k$, où $k\in \Rr$.}
    \good{$\displaystyle \int \frac{\mathrm{d}x}{x+1}=\ln (x+1)+k$, où $k\in \Rr$.}
    \bad{$\displaystyle \int \frac{\mathrm{d}x}{x+1}=\frac{-1}{(x+1)^2}+k$, où $k\in \Rr$.}
\end{answers}
\vskip3mm
\begin{explanations}
Avec $u=1+x$, on a : $\mathrm{d}u=\mathrm{d}x$,
$$\int \frac{\mathrm{d}x}{(x+1)^3}=\int u^{-3}\mathrm{d}u=\frac{-1}{2}u^{-2}+k=\frac{-1}{2(x+1)^2}+k,\; k\in \Rr,$$
et $\displaystyle \int \frac{\mathrm{d}x}{x+1}=\int \frac{\mathrm{d}u}{u}=\ln |u|+k=\ln (x+1)+k$, $k\in \Rr$.
\end{explanations}
\end{question}
}
